\chapter{System Design}
\label{chap:design}

This chapter describes the architecture of PROTEA at the level of design
decisions. Implementation details — specific libraries, migration tooling,
frontend components — are deferred to \cref{chap:implementation}.

% ────────────────────────────────────────────────────────────
\section{Requirements and Design Goals}
\label{sec:design-goals}

The design of PROTEA is governed by five requirements derived from the
limitations identified in \cref{chap:related}:

\begin{description}
  \item[R1 — Reproducibility.] A prediction produced today must be exactly
    reproducible in the future. This requires recording the ontology version,
    reference annotation set, and embedding model configuration used for every
    prediction run.
  \item[R2 — Scalability.] The system must handle reference sets of hundreds of
    thousands of proteins and query sets of thousands without holding all data
    in memory simultaneously.
  \item[R3 — Separation of concerns.] Domain logic (what to compute), execution
    flow (how jobs are dispatched and tracked), and infrastructure (database,
    message queue) must be independently replaceable.
  \item[R4 — Observability.] Every job must produce a structured audit trail so
    that failures can be diagnosed without replaying the computation.
  \item[R5 — Accessibility.] Researchers without machine-learning infrastructure
    expertise must be able to submit sequences and retrieve predictions through a
    web interface or a REST API.
\end{description}

% ────────────────────────────────────────────────────────────
\section{High-Level Architecture}
\label{sec:architecture}

PROTEA decomposes into four horizontal layers, illustrated in
\cref{fig:architecture}:

\begin{description}
  \item[Presentation layer.] A Next.js single-page application exposes the
    user-facing workflow: uploading FASTA files, triggering pipeline jobs,
    monitoring progress, and downloading results.
  \item[API layer.] A FastAPI application provides a REST interface. All state
    mutations flow through this layer; it writes job requests to PostgreSQL and
    publishes messages to RabbitMQ.
  \item[Worker layer.] A set of Python worker processes consume messages from
    RabbitMQ queues. Workers are stateless with respect to domain logic — they
    resolve and execute registered operations and update job state in the database.
  \item[Data layer.] PostgreSQL (extended with the pgvector extension for vector
    storage) is the single source of truth. RabbitMQ is used exclusively for
    task dispatch; it carries no persistent state.
\end{description}

\begin{figure}[htbp]
  \centering
  \resizebox{\textwidth}{!}{% ============================================================
%  Figure: PROTEA high-level architecture (four layers)
% ============================================================
\begin{tikzpicture}[
  node distance=6mm and 8mm,
  every node/.style={font=\small},
]

  % ── Presentation layer ──────────────────────────────────────
  \node[proteablock, minimum width=46mm] (web) {Next.js Web UI\\\footnotesize FASTA upload, job monitor, TSV export};

  % ── API layer ───────────────────────────────────────────────
  \node[proteablock, minimum width=46mm, below=of web] (api) {FastAPI REST API\\\footnotesize 17 routers · session injection};

  % ── Worker layer ────────────────────────────────────────────
  \node[proteablock, minimum width=22mm, below=14mm of api, xshift=-30mm] (qc) {QueueConsumer\\\footnotesize tracked jobs};
  \node[proteablock, minimum width=22mm, right=4mm of qc]                  (oc) {OperationConsumer\\\footnotesize batch sub-tasks};
  \node[proteablock, minimum width=20mm, right=4mm of oc]                  (gpu) {GPU Inference\\\footnotesize ESM-2 / ESM-C / ProstT5 / ProtT5 / Ankh};

  % ── Data layer ──────────────────────────────────────────────
  \node[proteastore, below=14mm of qc, xshift=10mm]    (pg)  {PostgreSQL 16\\\footnotesize + pgvector};
  \node[proteastore, right=14mm of pg]                 (mq)  {RabbitMQ\\\footnotesize 10 queues};

  % ── Layer brackets ──────────────────────────────────────────
  \begin{scope}[on background layer]
    \node[draw=proteagreen!40, very thick, dashed, rounded corners=4pt,
          fit=(web), inner sep=4mm, label={[font=\footnotesize\itshape, text=proteagreen!70!black]left:Presentation}] {};
    \node[draw=proteagreen!40, very thick, dashed, rounded corners=4pt,
          fit=(api), inner sep=4mm, label={[font=\footnotesize\itshape, text=proteagreen!70!black]left:API}] {};
    \node[draw=proteagreen!40, very thick, dashed, rounded corners=4pt,
          fit=(qc)(oc)(gpu), inner sep=3mm, label={[font=\footnotesize\itshape, text=proteagreen!70!black]left:Workers}] {};
    \node[draw=linkblue!40, very thick, dashed, rounded corners=4pt,
          fit=(pg)(mq), inner sep=3mm, label={[font=\footnotesize\itshape, text=linkblue!70!black]left:Data}] {};
  \end{scope}

  % ── Edges ───────────────────────────────────────────────────
  \draw[proteaarrow] (web) -- node[right, font=\footnotesize]{HTTPS / JSON} (api);
  \draw[proteaarrow] (api.south) -- ++(0,-3mm) -| (qc.north);
  \draw[proteaarrow] (api.south) -- ++(0,-3mm) -| (oc.north);
  \draw[proteaarrow] (api) -- ++(35mm,0) |- (mq.east);

  \draw[proteaarrow] (qc) -- (pg);
  \draw[proteaarrow] (oc) -- (pg);
  \draw[proteaarrow] (oc.south east) -- ++(0,-3mm) -| (gpu.south);
  \draw[proteaarrow] (qc.east) -- (oc.west);

  \draw[proteadashed] (mq.north) to[bend left=18] node[right, font=\footnotesize]{dispatch} (oc.east);
  \draw[proteadashed] (mq.north) to[bend left=10] (qc.east);

\end{tikzpicture}
}
  \caption[High-level architecture of PROTEA]{%
    High-level architecture of PROTEA. The four horizontal layers
    (presentation, API, workers, data) communicate exclusively through well-defined
    interfaces: HTTPS/JSON between the browser and the API, RabbitMQ for task
    dispatch, and PostgreSQL as the single source of truth. GPU inference is
    isolated in dedicated workers so that compute capacity can be scaled
    independently of API throughput.%
  }
  \label{fig:architecture}
\end{figure}

% ────────────────────────────────────────────────────────────
\section{The Operation Abstraction}
\label{sec:operation-abstraction}

The central design principle of PROTEA is the \textit{Operation protocol}: every
unit of domain logic is a class that implements exactly two members:

\begin{itemize}
  \item \texttt{name: str} — a unique string identifier used to route messages.
  \item \texttt{execute(session, payload, *, emit) -> OperationResult} — the
    computation itself.
\end{itemize}

Operations receive a SQLAlchemy session and a validated Pydantic payload; they
report progress through an \texttt{emit} callback that writes \texttt{JobEvent}
rows to the database. Operations do not manage their own sessions, do not
publish messages, and do not know about the queue infrastructure. This strict
interface makes them independently testable and trivially substitutable.

The \texttt{OperationRegistry} maps string names to instances at startup. Workers
resolve the correct operation by name when they receive a message, enabling
new operations to be added without modifying any worker code.

% ────────────────────────────────────────────────────────────
\section{Job Lifecycle and Worker Patterns}
\label{sec:job-lifecycle}

PROTEA distinguishes two worker patterns that correspond to different
observability requirements:

\subsection{QueueConsumer (tracked jobs)}
\label{subsec:queue-consumer}

Long-running, user-visible jobs — inserting proteins, loading ontologies,
computing embeddings, predicting GO terms — are submitted through the API as
\texttt{Job} rows with a JSONB payload. The workflow is:

\begin{enumerate}
  \item The API creates a \texttt{Job} row in state \texttt{QUEUED} and publishes
    the job UUID to the appropriate RabbitMQ queue.
  \item A \texttt{QueueConsumer} worker receives the UUID, opens a
    \textit{claim session}, transitions the job to \texttt{RUNNING}, and flushes
    the \texttt{started\_at} timestamp. The claim session is committed and closed
    before execution begins.
  \item A separate \textit{execute session} resolves the operation and calls
    \texttt{execute()}. On success the job transitions to \texttt{SUCCEEDED}; on
    exception to \texttt{FAILED}, with the error code and message stored on the
    row.
\end{enumerate}

The two-session design ensures that the claim is durable even if the execute
session rolls back. \texttt{JobEvent} rows, written via \texttt{emit()}, form an
append-only audit log that records progress milestones, error context, and
structured diagnostic fields.

\subsection{OperationConsumer (fire-and-forget sub-tasks)}
\label{subsec:operation-consumer}

GPU-intensive batch tasks — running the embedding model on a group of sequences,
or computing kNN predictions for a batch — do not create child \texttt{Job} rows.
Instead, the coordinator operation publishes \textit{operation messages}
(containing the payload directly, not a UUID) to dedicated queues.
\texttt{OperationConsumer} workers execute these messages and report progress
via an atomic increment on the parent job's counter, avoiding write contention
from hundreds of concurrent batches.

\Cref{fig:job-lifecycle} illustrates the two patterns and the queue topology.

\subsection{Parent-Child Job Hierarchy}
\label{subsec:parent-child}

Coordinator operations (\texttt{compute\_embeddings}, \texttt{predict\_go\_terms})
split work across many parallel workers using a parent-child pattern. The
coordinator returns \texttt{OperationResult(deferred=True)}, which signals
\texttt{BaseWorker} to \emph{not} transition the parent job to
\texttt{SUCCEEDED} immediately. Instead, the parent remains in state
\texttt{RUNNING} while batch workers process their messages independently.

The \texttt{Job} model carries two progress fields for this purpose:
\texttt{progress\_total} (set by the coordinator to the number of batches
dispatched) and \texttt{progress\_current} (incremented atomically by each
write worker). When the last write worker detects
$\texttt{progress\_current} = \texttt{progress\_total}$, it closes the parent
job as \texttt{SUCCEEDED}. This atomic counter design avoids write contention
from hundreds of concurrent batches updating the same row.

\subsection{RetryLaterError}
\label{subsec:retry-later}

When a shared resource is unavailable — for instance, the GPU is already
occupied by a running embedding job — an operation raises
\texttt{RetryLaterError} instead of failing permanently. \texttt{BaseWorker}
catches this exception and:

\begin{enumerate}
  \item Resets the job status back to \texttt{QUEUED}.
  \item Writes a \texttt{job.retry\_later} event recording the reason and
    the requested delay.
  \item Re-publishes the job UUID to its queue after the specified delay
    (typically 60 seconds).
\end{enumerate}

The embedding coordinator queue (\texttt{protea.embeddings}) is serialised
precisely because of this mechanism: at most one embedding coordinator runs at
a time, with subsequent requests queued and retried automatically.

\subsection{Cancellation}
\label{subsec:cancellation}

\texttt{POST /jobs/\{id\}/cancel} transitions a \texttt{QUEUED} or
\texttt{RUNNING} job to \texttt{CANCELLED}. If the job is already in a
terminal state (\texttt{SUCCEEDED}, \texttt{FAILED}, or \texttt{CANCELLED})
the request is a no-op. Any queued child jobs are also cancelled atomically
in the same transaction.

Cancellation of a \texttt{RUNNING} job is a \textit{soft cancel}: the database
row is marked \texttt{CANCELLED} immediately, but the worker process is not
interrupted. The worker will still complete its operation and attempt a final
status write; however, because \texttt{CANCELLED} is already committed, the
frontend reflects the cancelled state regardless of how the worker exits.

\begin{figure}[htbp]
  \centering
  \resizebox{\textwidth}{!}{% ============================================================
%  Figure: Job lifecycle and queue topology
%  Source: figures/fig_job_lifecycle.tex @ 57fa4d4
%  Static TikZ diagram. Verify fidelity against:
%    protea/api/app.py, protea/workers/base_worker.py,
%    scripts/manage.sh (10 queues: ping, jobs, training, embeddings,
%    embeddings.batch, embeddings.write, predictions, predictions.batch,
%    predictions.write, evaluations)
% ============================================================
\begin{tikzpicture}[
  node distance=6mm and 8mm,
  every node/.style={font=\small},
]

  % ── State machine (left column) ─────────────────────────────
  \node[proteastate]                       (queued)    {QUEUED};
  \node[proteastate, below=8mm of queued]  (running)   {RUNNING};
  \node[proteastate, below=8mm of running] (succeeded) {SUCCEEDED};
  \node[proteastate, right=10mm of succeeded] (failed)  {FAILED};
  \node[proteastate, left=10mm of succeeded]  (cancel) {CANCELLED};

  \draw[proteaarrow] (queued)  -- node[right, font=\footnotesize]{claim session} (running);
  \draw[proteaarrow] (running) -- node[right, font=\footnotesize]{execute session} (succeeded);
  \draw[proteaarrow] (running.east) -- ++(8mm,0) |- (failed.north);
  \draw[proteaarrow] (running.west) -- ++(-8mm,0) |- (cancel.north);

  % RetryLater self-loop
  \draw[proteadashed] (running.west) to[out=160, in=200, looseness=2.4]
    node[left, font=\footnotesize, align=right]{RetryLaterError\\(re-queue +60s)} (queued.west);

  \node[draw=black!20, dashed, rounded corners=3pt,
        fit=(queued)(running)(succeeded)(failed)(cancel),
        inner sep=4mm,
        label={[font=\footnotesize\itshape]above:Job state machine}] (sm) {};

  % ── Queue topology (right column) ───────────────────────────
  \node[proteablock, minimum width=30mm, right=24mm of queued.north east, anchor=north west] (coord) {Coordinator op\\\footnotesize \texttt{compute\_embeddings}\\\texttt{predict\_go\_terms}};

  \node[proteaqueue, below=8mm of coord]                      (qjobs)    {protea.jobs};
  \node[proteaqueue, below=2mm of qjobs]                      (qemb)     {protea.embeddings};
  \node[proteaqueue, below=2mm of qemb]                       (qebatch)  {protea.embeddings.batch};
  \node[proteaqueue, below=2mm of qebatch]                    (qewrite)  {protea.embeddings.write};
  \node[proteaqueue, below=2mm of qewrite]                    (qpbatch)  {protea.predictions.batch};
  \node[proteaqueue, below=2mm of qpbatch]                    (qpwrite)  {protea.predictions.write};
  \node[proteaqueue, below=2mm of qpwrite]                    (qping)    {protea.ping};
  \node[proteaqueue, below=2mm of qping]                      (qpred)    {protea.predictions};
  \node[proteaqueue, below=2mm of qpred]                      (qtrain)   {protea.training};
  \node[proteaqueue, below=2mm of qtrain]                     (qeval)    {protea.evaluations};

  \node[draw=BurntOrange!40, dashed, rounded corners=3pt,
        fit=(coord)(qjobs)(qeval),
        inner sep=4mm,
        label={[font=\footnotesize\itshape, text=BurntOrange!70!black]above:RabbitMQ topology (10 queues)}] (qg) {};

  % Coordinator → fan-out
  \draw[proteaarrow] (coord.south) -- ++(0,-1mm) -| (qebatch.north);
  \draw[proteaarrow] (coord.south) -- ++(0,-1mm) -| (qpbatch.north);

  % Batch → write chain
  \draw[proteaarrow] (qebatch.east) to[bend left=18] node[right, font=\footnotesize]{vectors} (qewrite.east);
  \draw[proteaarrow] (qpbatch.east) to[bend left=18] node[right, font=\footnotesize]{rows}    (qpwrite.east);

  % Atomic counter on parent job
  \node[draw=proteagreen!50, very thick, fill=proteagreen!8, rounded corners=3pt,
        right=12mm of qewrite, font=\footnotesize, align=center]
        (counter) {atomic\\progress\_current\\++};
  \draw[proteaarrow] (qewrite.east) -- (counter.west);
  \draw[proteaarrow] (qpwrite.east) to[bend right=18] (counter.south);
  \draw[proteadashed] (counter.north) to[bend left=12]
        node[right, font=\footnotesize]{closes parent} (running.east);

\end{tikzpicture}
}
  \caption[Job lifecycle and queue routing in PROTEA]{%
    Job lifecycle and queue routing in PROTEA. The state machine on the left
    is enforced by \texttt{BaseWorker} through two independent SQLAlchemy
    sessions: a short \emph{claim session} that commits the
    \texttt{QUEUED}~$\to$~\texttt{RUNNING} transition, and a longer
    \emph{execute session} that runs the operation and commits the terminal
    state. \texttt{RetryLaterError} resets the job to \texttt{QUEUED} and
    re-publishes it after a delay (typically 60\,s) so that GPU contention
    never produces hard failures. The right column shows the seven RabbitMQ
    queues; coordinator operations fan out work into the
    \texttt{*.batch}~$\to$~\texttt{*.write} chains, where every write worker
    atomically increments \texttt{progress\_current} on the parent job. The
    last write closes the parent as \texttt{SUCCEEDED}, avoiding write
    contention from hundreds of concurrent batches.%
  }
  \label{fig:job-lifecycle}
\end{figure}

% ────────────────────────────────────────────────────────────
\section{Queue Topology}
\label{sec:queue-topology}

PROTEA routes messages across seven RabbitMQ queues, each with a distinct
semantic purpose:

\begin{description}
  \item[\texttt{protea.ping}] Smoke-test queue. Used to verify end-to-end
    connectivity.
  \item[\texttt{protea.jobs}] General-purpose job queue. Handles protein
    insertion, metadata fetch, ontology loading, annotation loading, and the
    coordinator phase of embeddings and predictions.
  \item[\texttt{protea.embeddings}] Serialised embedding coordinator queue. At
    most one coordinator runs at a time; if the GPU is busy the coordinator
    raises \texttt{RetryLaterError} and is re-queued with a 60-second delay.
  \item[\texttt{protea.embeddings.batch}] GPU inference queue. Each message
    carries a batch of sequences; \texttt{OperationConsumer} workers run the
    forward pass of the protein language model and publish results to the
    write queue.
  \item[\texttt{protea.embeddings.write}] Bulk insert queue. Workers receive
    computed embedding vectors and write them to the \texttt{SequenceEmbedding}
    table via pgvector.
  \item[\texttt{protea.predictions.batch}] kNN prediction queue. Each message
    carries a batch of query embeddings; workers search the reference index and
    transfer GO annotations.
  \item[\texttt{protea.predictions.write}] Bulk prediction insert queue. Workers
    write \texttt{GOPrediction} rows for a batch of query--reference pairs.
\end{description}

This topology separates GPU-bound inference, I/O-bound database writes, and
CPU-bound coordination into independent queues that can be scaled independently
(e.g.\ \texttt{manage.sh scale protea.predictions.batch 4}).

% ────────────────────────────────────────────────────────────
\section{Data Model}
\label{sec:data-model}

The relational schema is organised into five logical groups:

\subsection{Sequence and Protein}
\label{subsec:dm-sequence}

\texttt{Sequence} stores deduplicated amino acid strings, keyed by MD5 hash.
\texttt{Protein} stores one row per UniProt accession; multiple accessions
(canonical and isoforms) may reference the same \texttt{Sequence}. This
deduplication prevents redundant embedding computation.

\subsection{Ontology and Annotations}
\label{subsec:dm-ontology}

\texttt{OntologySnapshot} records one complete GO release (OBO format), versioned
by the \texttt{obo\_version} string found in the OBO file header. \texttt{GOTerm}
and \texttt{GOTermRelationship} store the DAG within a snapshot.

\texttt{AnnotationSet} groups a batch of \texttt{ProteinGOAnnotation} rows by
source (\texttt{goa} or \texttt{quickgo}) and links them to an
\texttt{OntologySnapshot}. This design enables side-by-side comparison of
annotation sets from different sources or dates.

\subsection{Embeddings}
\label{subsec:dm-embeddings}

\texttt{EmbeddingConfig} is an immutable recipe: it records the model identifier,
chunking strategy, pooling method, and any other parameters that affect the
embedding geometry. Its UUID primary key is stable; changing any parameter
creates a new configuration.

\texttt{SequenceEmbedding} stores one pgvector \texttt{VECTOR} per
(\texttt{sequence}, \texttt{config}, \texttt{chunk}) triple. Chunking support
allows sequences that exceed the model's context window to be split into
overlapping segments.

\subsection{Predictions}
\label{subsec:dm-predictions}

\texttt{PredictionSet} is the result container, linking a query set, an embedding
configuration, an annotation set, and an ontology snapshot.
\texttt{GOPrediction} stores one row per (query protein, GO term) pair, including
the cosine distance from the nearest reference neighbour and the optional
feature-engineering columns described in \cref{sec:feature-engineering}.

\subsection{Reranker Models}
\label{subsec:dm-reranker}

\texttt{RerankerModel} records a trained LightGBM booster promoted from the
research sandbox described in \cref{sec:reranker-promotion}. The row carries
the artefact URI resolved by PROTEA's storage abstraction
(\texttt{file:///\ldots} for local filesystem or \texttt{s3://\ldots} for
MinIO), the \texttt{feature\_schema\_sha} fingerprint used to validate
feature-set alignment at inference time, the originating embedding
configuration and ontology snapshot, and the full experiment specification
YAML used to reproduce the run. The inline \texttt{model\_data} column is
retained only as a nullable compatibility slot for rows produced by the
pre-abstraction pipeline; new registrations never populate it.

\subsection{Query Sets}
\label{subsec:dm-querysets}

\texttt{QuerySet} represents a user-uploaded FASTA dataset submitted for custom
prediction queries. Each \texttt{QuerySetEntry} row preserves the original
accession header from the FASTA file and links to the deduplicated
\texttt{Sequence} row, reusing existing sequences when the amino-acid string is
already present in the database. Query sets are created via
\texttt{POST /query-sets} and can be referenced in subsequent embedding
computation and prediction jobs.

\subsection{Jobs}
\label{subsec:dm-jobs}

\texttt{Job} implements a state machine with states \texttt{QUEUED},
\texttt{RUNNING}, \texttt{SUCCEEDED}, \texttt{FAILED}, and \texttt{CANCELLED}.
\texttt{JobEvent} is an append-only log of timestamped, structured events
emitted during execution. Both tables use JSONB columns for extensible payloads
and metadata.

% ────────────────────────────────────────────────────────────
\section{Embedding Pipeline}
\label{sec:embedding-pipeline}

The embedding pipeline runs in three phases coordinated through the queue
topology described above:

\begin{enumerate}
  \item \textbf{Coordinator} (\texttt{compute\_embeddings}): queries the database
    for all sequences in the target set that do not yet have an embedding for the
    requested \texttt{EmbeddingConfig}; partitions them into fixed-size batches;
    publishes one operation message per batch to \texttt{protea.embeddings.batch}.
  \item \textbf{Batch inference} (\texttt{compute\_embeddings\_batch}): loads the
    protein language model on the configured device; tokenises the sequence
    batch; runs a forward pass; mean-pools the per-residue representations into
    a single vector per chunk; publishes the computed vectors to
    \texttt{protea.embeddings.write}.
  \item \textbf{Store} (\texttt{store\_embeddings}): receives the computed vectors
    and performs a bulk upsert into \texttt{SequenceEmbedding}.
\end{enumerate}

The coordinator is serialised (one at a time per the \texttt{protea.embeddings}
queue) to prevent concurrent model loads from exhausting GPU memory.
Batch and write workers can be scaled horizontally.

% ────────────────────────────────────────────────────────────
\section{Prediction Pipeline}
\label{sec:prediction-pipeline}

The prediction pipeline mirrors the structure of the embedding pipeline:

\begin{enumerate}
  \item \textbf{Coordinator} (\texttt{predict\_go\_terms}): loads the reference
    embeddings (all proteins in the annotation set that have embeddings) into
    a process-level cache compressed to float16; partitions the query set into
    batches; publishes one operation message per batch to
    \texttt{protea.predictions.batch}.
  \item \textbf{Batch prediction} (\texttt{predict\_go\_terms\_batch}): for each
    query, retrieves the $k$ nearest reference proteins using the configured
    search backend (NumPy exact cosine search or FAISS approximate search);
    transfers GO annotations from each neighbour; optionally computes pairwise
    alignment statistics (\gls{nw} and \gls{sw} via the parasail library) and
    taxonomic distance (via ete3 and the NCBI taxonomy); publishes the results
    to \texttt{protea.predictions.write}.
  \item \textbf{Store} (\texttt{store\_predictions}): bulk-inserts
    \texttt{GOPrediction} rows.
\end{enumerate}

The reference cache is held at the process level and is keyed by
(embedding config, annotation set). It holds at most one entry to bound memory
usage; a restart reloads from the database automatically.

When a job payload includes a \texttt{reranker\_model\_id}, the batch worker
performs one additional pass after the kNN stage: the booster referenced by
the registered \texttt{RerankerModel} is loaded through a sha-keyed on-disk
cache, the candidate rows are scored, and the resulting \texttt{reranker\_score}
values replace the kNN distance as the ranking key. The pass is skipped and a
\texttt{reranker.schema\_mismatch} event is emitted if the feature fingerprint
computed from the live pipeline does not match the one registered with the
booster; the job otherwise proceeds to the store stage unchanged. The
promotion mechanics are described in \cref{sec:reranker-promotion}.

% ────────────────────────────────────────────────────────────
\section{REST API Design}
\label{sec:api-design}

The API follows REST conventions. Resources are grouped into six routers:

\begin{description}
  \item[\texttt{/proteins}] Insert proteins from UniProt accession lists or FASTA.
  \item[\texttt{/annotations}] Load ontology snapshots and annotation sets.
  \item[\texttt{/embeddings}] Manage embedding configurations, trigger computation,
    submit prediction jobs, and export results.
  \item[\texttt{/query-sets}] Upload FASTA files and manage user query datasets.
  \item[\texttt{/jobs}] Track job state and stream structured events.
  \item[\texttt{/admin}] Administrative maintenance tasks.
\end{description}

\begin{sloppypar}
Session injection follows FastAPI's dependency system:
\texttt{app.state.session\_factory} is set at startup and injected into each
request handler, keeping router code free of infrastructure imports.
\end{sloppypar}
